\section{Theory}

We sometimes have to cite two scientific publications or two reports, or just two citations together in general. For that, you can use \verb|\cite{citation1, citation2}|, i.e., separate the citations with a comma, like~\parencite{Beck2000application, Beck2005timeseries, Milano2018}.

When you are closely following a book to explain something, which is often the case in a theory section, you can write at the start of the section you are about to introduce: ``This section follows closely the reference book by Sumiyoshi Abe and Yuko Okamoto, \emph{Nonextensive Statistical Mechanics and Its Applications}~\parencite[p.~32]{Abe2001nonextensive}.''

A Master's~\cite{LastName2045norwegian} or a PhD thesis~\cite{Temult2038binding} should include the name of the university wherein it was written, as well as the year. Moreover, it should include a URL to the work, when available.

Articles and proceedings are practically the same in science, but in your bibliography, you should use \verb|@article| for articles and \verb|@inproceedings| for proceedings. The work by Jozefowicz et al. (2015)~\cite{Jozefowicz2015empirical} is an example of a proceedings.
